% GENERATED FILE. Edit canonical lessons, then run npm run generate:books.
% canonical-source-hash: fnv1a64:e56a4282862265e0
% canonical-lessons: JA-C314-sukkari, JA-C314-itsunomanika, JA-C314-yagate, JA-C314-heru, JA-C314-omoitsuku

\chapter{Everything Has Changed}
\label{ch:ja-everything-has-changed}
\hlchaptermodality{314}

\begin{chapteropening}
\textbf{By the end of this chapter:} \emph{I can say how a place has changed over time.}

Complete the last lesson of chapter 314: I can say how a place has changed over time.
\end{chapteropening}

\section[\texorpdfstring{\ja{すっかり} (sukkari)}{すっかり (sukkari)}]{\ja{すっかり} (sukkari) — completely; quite}

\label{lesson:JA-C314-sukkari}



\begin{quote}
Before the new one: say the Japanese for the other day, then the Japanese for ruins.
\end{quote}

\subsection*{You'll want to know: \ja{すっかり}}
\textbf{\ja{すっかり}} — \emph{sukkari} — \textquotedblleft{}completely; quite\textquotedblright{}.

\begin{quote}
\textbf{\ja{まちは} \ja{すっかり} \ja{かわりました。}} — \emph{machi wa sukkari kawarimashita} — \textquotedblleft{}The town has changed completely.\textquotedblright{}
\end{quote}

\begin{cousinweb}
\textbf{\ja{すっかり}} (\emph{sukkari}) is a native Japanese mimetic word, inherited rather than borrowed from Chinese, and its older origin is not settled (JMdict).
\end{cousinweb}

\begin{grammarlens}[title={What you've built}]
One more word for telling what happened.
\end{grammarlens}

\subsection*{Your turn}
\begin{itemize}
\raggedright
  \item \emph{Say it:} \emph{sukkari}
  \item \emph{Say it:} \emph{sukkari}, once more
  \item \emph{Say it:} \emph{machi wa sukkari kawarimashita} — The town has changed completely.
  \item \emph{Recall:} say the Japanese for the other day, then the Japanese for ruins, then say \emph{sukkari} again
\end{itemize}

\begin{tcolorbox}[breakable,colback=teal!4,colframe=teal!35!black,title={Before you move on}]
What does \ja{すっかり} mean? (\textquotedblleft{}Completely; quite\textquotedblright{}.) Say it once more.
\end{tcolorbox}
\section[\texorpdfstring{\ja{いつのまにか} (itsunomanika)}{いつのまにか (itsunomanika)}]{\ja{いつのまにか} (itsunomanika) — before one knows it}

\label{lesson:JA-C314-itsunomanika}



\begin{quote}
Before the new one: say the Japanese for ruins, then the Japanese for completely.
\end{quote}

\subsection*{You'll want to know: \ja{いつのまにか}}
\textbf{\ja{いつのまにか}} — \emph{itsunomanika} — \textquotedblleft{}before one knows it\textquotedblright{}.

\begin{quote}
\textbf{\ja{いつのまにか} \ja{ゆうがたに} \ja{なっていました。}} — \emph{itsunomanika yūgata ni natte imashita} — \textquotedblleft{}Before we knew it, evening had come.\textquotedblright{}
\end{quote}

\begin{cousinweb}
\textbf{\ja{いつのまにか}} (\emph{itsunomanika}) is a native Japanese word (wago), inherited rather than borrowed from Chinese, built from \emph{itsu}, \textquotedblleft{}when\textquotedblright{}, and \emph{ma}, \textquotedblleft{}an interval\textquotedblright{} (JMdict, KANJIDIC2).
\end{cousinweb}

\begin{grammarlens}[title={What you've built}]
One more word for telling what happened.
\end{grammarlens}

\subsection*{Your turn}
\begin{itemize}
\raggedright
  \item \emph{Say it:} \emph{itsunomanika}
  \item \emph{Say it:} \emph{itsunomanika}, once more
  \item \emph{Say it:} \emph{itsunomanika yūgata ni natte imashita} — Before we knew it, evening had come.
  \item \emph{Recall:} say the Japanese for ruins, then the Japanese for completely, then say \emph{itsunomanika} again
\end{itemize}

\begin{tcolorbox}[breakable,colback=teal!4,colframe=teal!35!black,title={Before you move on}]
What does \ja{いつのまにか} mean? (\textquotedblleft{}Before one knows it\textquotedblright{}.) Say it once more.
\end{tcolorbox}
\section[\texorpdfstring{\ja{やがて} (yagate)}{やがて (yagate)}]{\ja{やがて} (yagate) — before long; in time}

\label{lesson:JA-C314-yagate}



\begin{quote}
Before the new one: say the Japanese for completely, then the Japanese for before one knows it.
\end{quote}

\subsection*{You'll want to know: \ja{やがて}}
\textbf{\ja{やがて}} — \emph{yagate} — \textquotedblleft{}before long; in time\textquotedblright{}.

\begin{quote}
\textbf{\ja{やがて} \ja{あめが} \ja{やみました。}} — \emph{yagate ame ga yamimashita} — \textquotedblleft{}Before long the rain stopped.\textquotedblright{}
\end{quote}

\begin{cousinweb}
\textbf{\ja{やがて}} (\emph{yagate}) is a native Japanese word (wago), inherited rather than borrowed from Chinese, and it is usually written in kana (JMdict).
\end{cousinweb}

\begin{grammarlens}[title={What you've built}]
One more word for telling what happened.
\end{grammarlens}

\subsection*{Your turn}
\begin{itemize}
\raggedright
  \item \emph{Say it:} \emph{yagate}
  \item \emph{Say it:} \emph{yagate}, once more
  \item \emph{Say it:} \emph{yagate ame ga yamimashita} — Before long the rain stopped.
  \item \emph{Recall:} say the Japanese for completely, then the Japanese for before one knows it, then say \emph{yagate} again
\end{itemize}

\begin{tcolorbox}[breakable,colback=teal!4,colframe=teal!35!black,title={Before you move on}]
What does \ja{やがて} mean? (\textquotedblleft{}Before long; in time\textquotedblright{}.) Say it once more.
\end{tcolorbox}
\section[\texorpdfstring{\ja{へる} (heru)}{へる (heru)}]{\ja{へる} (heru) — to decrease; to go down}

\label{lesson:JA-C314-heru}



\begin{quote}
Before the new one: say the Japanese for before one knows it, then the Japanese for before long.
\end{quote}

\subsection*{You'll want to know: \ja{へる}}
\textbf{\ja{へる}} — \emph{heru} — \textquotedblleft{}to decrease; to go down\textquotedblright{}.

\begin{quote}
\textbf{\ja{いなかの} \ja{ひとが} \ja{へっています。}} — \emph{inaka no hito ga hette imasu} — \textquotedblleft{}The number of people in the countryside is falling.\textquotedblright{}
\end{quote}

Polite \textbf{\ja{へります}}, \emph{herimasu}; a godan verb. Past (polite): \textbf{\ja{むらの} \ja{ひとが} \ja{へりました。}}, \emph{mura no hito ga herimashita}.

\begin{cousinweb}
\textbf{\ja{へる}} (\emph{heru}) is a native Japanese word (wago), inherited rather than borrowed from Chinese; its kanji was attached later as a kun'yomi (JMdict).
\end{cousinweb}

\begin{grammarlens}[title={What you've built}]
One more word for telling what happened.
\end{grammarlens}

\subsection*{Your turn}
\begin{itemize}
\raggedright
  \item \emph{Say it:} \emph{heru}
  \item \emph{Say it:} \emph{heru}, once more
  \item \emph{Say it:} \emph{inaka no hito ga hette imasu} — The number of people in the countryside is falling.
  \item \emph{Recall:} say the Japanese for before one knows it, then the Japanese for before long, then say \emph{heru} again
\end{itemize}

\begin{tcolorbox}[breakable,colback=teal!4,colframe=teal!35!black,title={Before you move on}]
What does \ja{へる} mean? (\textquotedblleft{}To decrease; to go down\textquotedblright{}.) Say it once more.
\end{tcolorbox}
\section[\texorpdfstring{\ja{おもいつく} (omoitsuku)}{おもいつく (omoitsuku)}]{\ja{おもいつく} (omoitsuku) — to think of; to hit upon}

\label{lesson:JA-C314-omoitsuku}



\begin{quote}
Before the new one: say the Japanese for before long, then the Japanese for to decrease.
\end{quote}

\subsection*{You'll want to know: \ja{おもいつく}}
\textbf{\ja{おもいつく}} — \emph{omoitsuku} — \textquotedblleft{}to think of; to hit upon\textquotedblright{}.

\begin{quote}
\textbf{\ja{いい} \ja{ほうほうを} \ja{おもいつきました。}} — \emph{ii hōhō o omoitsukimashita} — \textquotedblleft{}I have had a good idea how to do it.\textquotedblright{}
\end{quote}

Polite \textbf{\ja{おもいつきます}}, \emph{omoitsukimasu}; a godan verb. Past (polite): \textbf{\ja{ちちが} \ja{いい} \ja{ほうほうを} \ja{おもいつきました。}}, \emph{chichi ga ii hōhō o omoitsukimashita}.

\begin{cousinweb}
\textbf{\ja{おもいつく}} (\emph{omoitsuku}) is a native Japanese word (wago), inherited rather than borrowed from Chinese, built from \emph{omoi}, \textquotedblleft{}thinking\textquotedblright{} (from \emph{omou}), and \emph{tsuku}, \textquotedblleft{}to arrive at\textquotedblright{} (JMdict, KANJIDIC2).
\end{cousinweb}

\begin{grammarlens}[title={What you've built}]
One more word for telling what happened.
\end{grammarlens}

\subsection*{Your turn}
\begin{itemize}
\raggedright
  \item \emph{Say it:} \emph{omoitsuku}
  \item \emph{Say it:} \emph{omoitsuku}, once more
  \item \emph{Say it:} \emph{ii hōhō o omoitsukimashita} — I have had a good idea how to do it.
  \item \emph{Recall:} say the Japanese for before long, then the Japanese for to decrease, then say \emph{omoitsuku} again
\end{itemize}

\begin{tcolorbox}[breakable,colback=teal!4,colframe=teal!35!black,title={Before you move on}]
What does \ja{おもいつく} mean? (\textquotedblleft{}To think of; to hit upon\textquotedblright{}.) Say it once more.
\end{tcolorbox}
